%! Orthogonality of the Four Subspaces — Unit 4, Lesson 4.1 — Warm-Up
\documentclass[10pt]{article}
\usepackage{linalg-article}
\usepackage{linalg-boxes}
\newcommand{\TallMath}[1]{$\displaystyle #1\rule[-1.4em]{0pt}{3.2em}$}
\newcommand{\vv}{\mathbf{v}}
\newcommand{\ww}{\mathbf{w}}
\newcommand{\xx}{\mathbf{x}}
\newcommand{\T}{\mathsf{T}}

\begin{document}

\pageheader{Unit 4, Lesson 4.1}{Warm-Up}
\namedateperiod

\vspace{0.12in}

\begin{notesbox}{Getting Ready: Skills We'll Use Today}
\small These three ideas (Lessons 1.2, 3.2, 3.5) set up today's lesson. Answer quickly.

\vspace{4pt}
\textbf{1. Perpendicular test.} Are $\vv=\begin{bmatrix} 2 \\ 1 \\ 3 \end{bmatrix}$ and
$\ww=\begin{bmatrix} 1 \\ -5 \\ 1 \end{bmatrix}$ perpendicular? Compute the dot product and decide.
\[
  \vv\cdot\ww = \blank{3.4cm}
\]

\vspace{4pt}
\textbf{2. A nullspace direction.} For $A=\begin{bmatrix} 1 & 2 \\ 2 & 4 \end{bmatrix}$, the equation
$A\xx=\mathbf{0}$ reduces to $x_1+2x_2=0$. Give one nonzero vector $\xx$ in the nullspace.
\[
  \xx = \begin{bmatrix}\rule{0pt}{1.1em}\blank{1.1cm}\\[2pt]\blank{1.1cm}\end{bmatrix}
\]

\vspace{4pt}
\textbf{3. Counting dimensions (recall 3.5).} A matrix is $3\times 3$ with rank $r=2$. Fill in the
dimensions of its four subspaces:
\[
  \dim C(A^{\T}) = \blank{1.0cm},\quad \dim N(A) = \blank{1.0cm},\quad
  \dim C(A) = \blank{1.0cm},\quad \dim N(A^{\T}) = \blank{1.0cm}.
\]
\end{notesbox}

\end{document}
